\documentclass{article}
\usepackage[UTF8]{ctex}
\usepackage{amsmath,amssymb,amsfonts}
\usepackage{enumitem}
\usepackage{geometry}
\geometry{a4paper, margin=1in}
% === 错题本排版优化（不参与题目标识） ===
\setlength{\parindent}{0pt}
\setlength{\parskip}{0.35em plus 0.10em minus 0.08em}
\setlength{\emergencystretch}{3em}
\setlist{itemsep=0.25em, topsep=0.40em, parsep=0pt}
\allowdisplaybreaks[2]
% === 错题本排版优化结束 ===
\begin{document}
\section*{上海大学2025年高等代数真题}
\section*{39．上海大学2025年高等代数试卷}
一．填空题．每题5分，共25分．\\

\subsection{题目1}
1．设 \(p, p_{1}, p_{2}, \cdots, p_{s}\) 是互不相同的素数，则多项式 \(f(x)=x^{p}-p_{1} p_{2} \cdots p_{s}\) 在有理数域上 \(\_\_\_\_\)（可约／不可约）。

\subsection{题目2}
2．（可能有误）三阶矩阵 \(A\) 的行列式因子为 \(1, x-2,(x-2)^{2}(x-3)^{2}\) ，则 \(A\) 的 Jordan 标准形为\\
\(\_\_\_\_\)。

\subsection{题目3}
3．设 5 阶矩阵 \(A\) 的各行元素之和为 0 ，且 \(r(A)=4\) ，则 \(A X=0\) 的通解为 \(\_\_\_\_\)．

\subsection{题目4}
4．线性空间 \(V\) 是由3阶实矩阵关于矩阵加法和数乘运算构成的实线性空间，记

\[
V_{1}=\left\{A \in V \mid A^{\mathrm{T}}=-A\right\}, V_{2}=\left\{\left(a_{i j}\right)_{3 \times 3} \mid a_{i j}=0, \forall i<j\right\} .
\]

则 \(\operatorname{dim}\left(V_{1}+V_{2}\right)+\operatorname{dim}\left(V_{1} \cap V_{2}\right)=\) \(\_\_\_\_\)．

\subsection{题目5}
5．若二次型 \(f\left(x_{1}, x_{2}, x_{3}\right)=a x_{1}^{2}-a x_{2}^{2}+(a-1) x_{3}^{2}+2 x_{1} x_{3}-2 x_{2} x_{3}\) ，则 \(a=\) \(\_\_\_\_\)时，\(f\) 在 \(\mathbb{R}\) 上的规范形为 \(y_{1}^{2}-y_{2}^{2}\) 。

二．判断题．每题5分，共25分．正确请说明理由，错误请举出反例．\\

\subsection{题目6}
6．多项式 \(x^{3}-1\) 与 \(x^{5}-1\) 的最大公因式为1。

\subsection{题目7}
7．若向量组 \(\alpha_{1}, \alpha_{2}, \alpha_{3}\) 线性无关，则 \(\alpha_{1}+\alpha_{2}, \alpha_{2}-\alpha_{3}, \alpha_{1}+2 \alpha_{3}\) 也线性无关。

\subsection{题目8}
8．若 \(U, V\) 时 \(\mathbb{F}^{n}\) 的子空间，且 \(\operatorname{dim} U+\operatorname{dim} V=n\) ，则 \(\mathbb{F}^{n}=U \oplus V\) 。

\subsection{题目9}
9．对任意实矩阵 \(A_{m \times n}\) ，都有 \(E+A^{\mathrm{T}} A\) 为正定矩阵。

\subsection{题目10}
10．设 \(U\) 是实内积空间 \(V\) 上的一个线性变换的一个有限维不变子空间，则 \(U\) 的正交补空间 \(U^{\perp}\) 也是线性变换的不变子空间。
\end{document}
